% Ejemplo: dos límites calculados con la definición (castellano).
%
% Dos frames, dos diapositivas. En los apuntes salen seguidos, cada uno con su
% título como subapartado.

\begin{frame}
\didactatitle{Un límite por la definición}

\begin{example}
Probemos que $\lim_{x \to 2} (3x - 1) = 5$. Dado $\varepsilon > 0$, buscamos
$\delta > 0$ tal que $0 < |x - 2| < \delta$ implique
$|(3x - 1) - 5| < \varepsilon$. Como
\[
  |(3x - 1) - 5| = |3x - 6| = 3\,|x - 2| ,
\]
basta tomar $\delta = \varepsilon / 3$: si $0 < |x - 2| < \varepsilon/3$,
entonces $|(3x - 1) - 5| = 3\,|x - 2| < \varepsilon$.
\end{example}

\onlynotes{%
  El método es siempre el mismo: escribir $|f(x) - L|$ en términos de
  $|x - a|$ y leer de ahí cuánto hay que acercarse.
}
\end{frame}

\begin{frame}
\didactatitle{Acotar antes de elegir}

\begin{example}
Probemos que $\lim_{x \to 1} x^2 = 1$. Ahora
$|x^2 - 1| = |x - 1|\,|x + 1|$, y el factor $|x + 1|$ depende de $x$. Lo acotamos: si $|x - 1| < 1$, entonces
$0 < x < 2$ y $|x + 1| < 3$, así que $|x^2 - 1| < 3\,|x - 1|$.

\dpause

Tomamos
\begin{keyformula}
  \delta = \min\{1, \varepsilon/3\} .
\end{keyformula}
Si $0 < |x - 1| < \delta$, se cumplen a la vez $|x - 1| < 1$ y
$|x - 1| < \varepsilon/3$, luego $|x^2 - 1| < 3\,|x - 1| < \varepsilon$.
\end{example}

\begin{commonmistake}
Toman $\delta = \varepsilon / |x + 1|$. Pero $\delta$ se elige antes que $x$:
puede depender de $\varepsilon$ y de $a$, nunca de $x$.
\end{commonmistake}
\end{frame}
